problem:  
Let \(G\) be a step‑\(s\ge2\) Carnot group with Carnot–Carathéodory distance \(d_{CC}\) and Haar measure \(\mathfrak m\). Suppose \(E\subset G\) has **locally finite horizontal perimeter**, and that at \(|D_H\chi_E|\)-a.e. point \(x\in\partial^*E\) (the reduced boundary), every measure‑theoretic blowup \(\delta_{x,r}(E):=\delta_{1/r}(x^{-1}E)\) converges in \(L^1_{\mathrm{loc}}\) to a vertical half‑space \(H_x\), with horizontal normal \(\nu_E(x)\in V_1\).

Assume a **quantitative tangent-flatness (excess‑decay) condition**: there exist constants \(\beta>0\) and \(C>0\) such that for all sufficiently small \(r>0\),
\[
\mathcal E(x,r) := \frac{1}{r^{Q-1}} \int_{B(x,r)} 
|\nu_E(y) - \nu_E(x)|^2 \, d|D_H\chi_E|(y)
\ \le C r^{2\beta}.
\]
That is, the horizontal normal oscillates in \(L^2\) with Hölder exponent \(\beta\).

**Question:**  
Does this quantitative tangent-flatness imply that, in some neighborhood of \(x\), the reduced boundary \(\partial^*E\) coincides (up to negligible sets) with an **intrinsic \(C^{1,\beta'}\) hypersurface** for some \(\beta'>0\)?  
In particular, can one deduce intrinsic Lipschitz graph structure or quantitative rectifiability of \(\partial^*E\) under this excess‑decay assumption in arbitrary step‑\(s\) Carnot groups?

Why is it a "good" problem:  
This problem sits precisely at a research frontier linking **geometric measure theory**, **sub‑Riemannian geometry**, and **quantitative analysis of perimeter measures**. In Euclidean spaces, De Giorgi’s excess‑decay lemma implies that Hölder control on the oscillation of unit normals forces \(C^{1,\alpha}\) regularity of boundaries. In Carnot groups, despite the known existence of half‑space tangents and intrinsic Lipschitz classifications in special regimes (e.g. constant or continuous normals), no general **quantitative sub‑Riemannian De Giorgi theory** is known.  
The question asks whether a purely metric–measure assumption on the decay of horizontal normal oscillations suffices to induce intrinsic smoothness—an extension requiring new analytic and geometric tools adapted to noncommutativity, anisotropic dilations, and lack of a linear projection structure.  
A positive answer would generalize the Euclidean excess‑decay principle into the sub‑Riemannian world; a negative one would expose deep structural obstructions to local regularity in higher‑step groups. The problem is simply stated but conceptually profound, creating a bridge between **quantitative blowup theory** and **geometric regularity** in non‑Euclidean metric geometries.